
\chapter{Sistemas de Sinalização}
\label{cap:09}

A vida multicelular exige coordenação. Uma célula do fígado não pode simplesmente
decidir proliferar ou morrer em isolamento --- ela recebe, interpreta e integra
sinais de vizinhos, de órgãos distantes e do ambiente externo antes de agir. Essa
capacidade de detecção, transmissão e resposta constitui o que chamamos de
\emph{sinalização celular}, tema central deste capítulo.

O desafio é duplo. Do lado biológico, interessa entender como moléculas
extracelulares, muitas delas incapazes de atravessar a membrana plasmática,
conseguem alterar o estado do núcleo e do metabolismo. Do lado quantitativo,
interessa compreender por que cascatas de fosforilação exibem comportamentos
qualitatividade distintos --- respostas em degrau, oscilações, memória --- que
não estariam presentes em reações isoladas. A sinalização celular é, portanto,
um dos objetos de estudo mais ricos da biologia de sistemas: é um texto escrito
em moléculas, mas lido em propriedades emergentes.

\section{Fundamentos da Comunicação Celular}
\label{sec:09-fundamentos}

\begin{definicaoenv}[Sinalização Celular]
Sinalização celular é o processo pelo qual células detectam, transmitem e
respondem a sinais químicos ou físicos, permitindo comunicação intra e
intercelular. O processo envolve quatro etapas: emissão do sinal, recepção por
receptor específico, transdução intracelular e produção de resposta.
\end{definicaoenv}

A lógica geral da sinalização pode ser resumida como: \textbf{ligante}
$\rightarrow$ \textbf{receptor} $\rightarrow$ \textbf{transdução} $\rightarrow$
\textbf{resposta}. Em cada etapa, o sinal pode ser amplificado, modificado,
integrado com outros sinais ou extinto. Essa cadeia não é uma correia de
transmissão passiva --- cada componente é um processador que contribui com
lógica própria ao resultado final.

\subsection{Classificação por Alcance}
\label{subsec:09-alcance}

A primeira distinção relevante para sistemas de sinalização é o alcance
espacial do sinal. Sinais \textbf{autócrinos} são secretados e reconhecidos pela
própria célula emissora --- padrão frequente em cânceres, onde células neoplásicas
produzem os fatores de crescimento que estimulam sua própria proliferação.
Sinais \textbf{parácrinos} atuam sobre células vizinhas a curta distância: os
neurotransmissores liberados na fenda sináptica percorrem menos de 40 nanômetros
antes de atingir os receptores pós-sinápticos. Sinais \textbf{endócrinos} ---
hormônios como a insulina, os esteroides adrenais e a tiroxina --- são lançados
na corrente sanguínea e podem atingir alvos em qualquer órgão do organismo,
percorrendo metros em minutos. Há ainda a sinalização \textbf{justácrina}, que
exige contato físico direto entre membranas, como na via Notch-Delta, onde o
ligante e o receptor residem em células adjacentes.

Cada modalidade impõe compromissos distintos entre especificidade, velocidade e
alcance. A sinalização parácrina é rápida e localizada, mas de alcance limitado.
A endócrina atinge todos os tecidos, mas requer um sistema de distribuição
dedicado e é inevitavelmente lenta. A justácrina garante especificidade absoluta
de célula a célula, ao custo de depender de contato físico.

\section{Receptores Celulares e Mecanismos de Ativação}
\label{sec:09-receptores}

O ponto de entrada de qualquer sinal é o receptor. Receptores são proteínas que
reconhecem um ligante com alta especificidade e, ao fazê-lo, sofrem uma mudança
conformacional que inicia a cascata intracelular. Podem estar localizados na
membrana plasmática, no citoplasma ou no núcleo, dependendo das propriedades
físico-químicas do ligante.

A classificação clássica distingue cinco famílias: (i) receptores acoplados à
proteína G (GPCRs), (ii) receptores tirosina quinase (RTKs), (iii) receptores de
canais iônicos, (iv) receptores nucleares e (v) receptores de citocinas
(JAK-STAT). Cada família utiliza mecanismos de transdução distintos, com
diferentes velocidades de resposta e repertórios de efeitos.

\subsection{Receptores Acoplados à Proteína G (GPCRs)}
\label{subsec:09-gpcr}

Os GPCRs são a maior família de receptores em humanos, com aproximadamente 800
membros, e constituem o alvo de cerca de 50\% dos fármacos aprovados clinicamente.
Sua estrutura característica é a serpentina de sete domínios transmembrana que
atravessa a bicamada lipídica em zigue-zague. O domínio extracelular reconhece
o ligante; o domínio intracelular interage com a proteína G heterotrimérica
($\alpha\beta\gamma$) no estado inativo, ligada ao GDP.

O ciclo de ativação é elegante. Quando o ligante se une ao GPCR, induz uma
mudança conformacional que facilita a troca de GDP por GTP na subunidade
$G\alpha$. O complexo $G\alpha$-GTP se dissocia de $G\beta\gamma$, e ambas as
partes --- $G\alpha$-GTP e o dímero $G\beta\gamma$ --- podem ativar efetores
independentemente. A terminação ocorre por hidrólise intrínseca do GTP a GDP
pela atividade GTPase de $G\alpha$, restaurando o estado inativo e reciclando o
receptor.

Os subtipos de proteína G determinam a identidade do sinal downstream. $G_s$
estimula a adenilil ciclase, elevando os níveis intracelulares de cAMP; $G_i$
inibe a mesma enzima; $G_q$ ativa a fosfolipase C$\beta$, que cliva PIP$_2$ em
IP$_3$ e DAG, mobilizando Ca$^{2+}$ e ativando a proteína quinase C; $G_{12/13}$
regula a organização do citoesqueleto via RhoGEFs. A mesma lógica de acoplamento
--- ligante $\to$ conformação $\to$ seleção de via --- pode assim gerar respostas
metabólicas completamente distintas dependendo do subtipo de G acoplado ao receptor.

\subsection{Cinética de Ligação Ligante-Receptor}
\label{subsec:09-cinetica}

A interação entre um ligante L e seu receptor R é descrita no equilíbrio pelo
modelo de ligação simples:
\begin{equation}
L + R \underset{k_{\mathrm{off}}}{\overset{k_{\mathrm{on}}}{\rightleftharpoons}} LR
\end{equation}

A constante de dissociação no equilíbrio é definida como:
\begin{equation}
K_d = \frac{k_{\mathrm{off}}}{k_{\mathrm{on}}} = \frac{[L][R]}{[LR]}
\end{equation}

A fração de receptores ocupados, $\theta$, segue a equação de Hill para $n = 1$
(hipérbole de Michaelis):
\begin{equation}
\theta = \frac{[LR]}{[R]_{\mathrm{total}}} = \frac{[L]}{K_d + [L]}
\end{equation}

O $K_d$ tem interpretação direta: é a concentração de ligante que ocupa metade
dos receptores. Quanto menor o $K_d$, maior a afinidade. Receptores de alta
afinidade como os de insulina têm $K_d$ na faixa de 0.1--1 nM; receptores de
neuro\-transmissores, na faixa de $\mu$M. O código abaixo implementa a curva de
ligação e permite visualizar o efeito do $K_d$ na resposta:

\begin{lstlisting}[language=Python, caption=Curva de ligacao ligante-receptor]
import numpy as np
import matplotlib.pyplot as plt

def receptor_binding(L, Kd, Bmax):
    # Fracao de receptores ocupados (modelo de equilibrio)
    return (Bmax * L) / (Kd + L)

Kd = 10    # nM
Bmax = 100 # unidades arbitrarias

L = np.logspace(-1, 3, 200)  # 0.1 a 1000 nM
bound = receptor_binding(L, Kd, Bmax)

plt.figure(figsize=(8, 5))
plt.semilogx(L, bound, 'b-', linewidth=2, label=f'Kd = {Kd} nM')
plt.axhline(Bmax/2, color='r', linestyle='--', label='50% ocupacao')
plt.axvline(Kd, color='r', linestyle='--')
plt.xlabel('[Ligante] (nM)')
plt.ylabel('Receptores Ocupados')
plt.title('Ligacao Ligante-Receptor (equilibrio)')
plt.legend()
plt.grid(True, alpha=0.3)
plt.tight_layout()
plt.savefig('ligacao_receptor.png', dpi=150)
\end{lstlisting}

Quando o receptor é capaz de assumir três estados --- livre (R), ligado (LR) e
ativo ($R^*$) --- o modelo se torna um sistema de três equações diferenciais
ordinárias:
\begin{align}
\frac{d[R]}{dt} &= -k_{\mathrm{on}}[L][R] + k_{\mathrm{off}}[LR]
                   + k_{\mathrm{inact}}[R^*] \\
\frac{d[LR]}{dt} &= k_{\mathrm{on}}[L][R] - k_{\mathrm{off}}[LR]
                    - k_{\mathrm{act}}[LR] \\
\frac{d[R^*]}{dt} &= k_{\mathrm{act}}[LR] - k_{\mathrm{inact}}[R^*]
\end{align}

No estado estacionário ($d/dt = 0$), a concentração de receptor ativo é:
\begin{equation}
[R^*]_{ss} = \frac{k_{\mathrm{act}} \cdot k_{\mathrm{on}}[L][R]_{\mathrm{total}}}
{k_{\mathrm{inact}}(k_{\mathrm{off}} + k_{\mathrm{act}})
 + k_{\mathrm{act}} \cdot k_{\mathrm{on}}[L]}
\end{equation}

Essa expressão é análoga à equação de Michaelis-Menten: o numerador cresce
linearmente com o estímulo, o denominador satura, e a resposta tem a forma de
hipérbole retangular.


\subsection{Receptores Tirosina Quinase (RTKs)}
\label{subsec:09-rtk}

Os RTKs constituem a segunda grande família de receptores membranares. Caracterizam-se
por um domínio extracelular de ligação ao ligante, um domínio transmembrana único e
um domínio citoplasmático com atividade tirosina quinase intrínseca. Diferentemente
dos GPCRs, os RTKs são ativados por \textbf{dimerização}: a ligação do ligante
induz a aproximação de dois monômeros, e a proximidade de seus domínios quinase
permite a \emph{autofosforilação} cruzada em resíduos de tirosina. Essas
tirosinas fosforiladas recrutam diretamente proteínas adaptadoras contendo
domínios SH2 ou PTB, iniciando cascatas a jusante.

As principais famílias de RTKs incluem o receptor de EGF (EGFR/ErbB), o receptor
de PDGF (PDGFR), os receptores de FGF (FGFR1--4), o receptor de insulina (InsR)
e o receptor de VEGF (VEGFR). Cada família é especializada em diferentes
processos biológicos: EGFR dirige proliferação e diferenciação epitelial;
VEGFR medeia angiogênese; InsR coordena resposta metabólica à glicose. A
desregulação dos RTKs --- por superexpressão, mutação ativadora ou amplificação
gênica --- é uma das marcas do câncer: o glioblastoma tem amplificação de EGFR;
o câncer colorretal frequentemente superexpressa HER2 (um ErbB); os inibidores
de tirosina quinase (TKIs) como gefitinib e erlotinib foram desenvolvidos
justamente para silenciar esses receptores constitutivamente ativos.

\subsection{Receptores de Canais Iônicos e Nucleares}
\label{subsec:09-canais-nucleares}

Há duas famílias com localização e mecanismo distintos. Os \textbf{receptores de
canais iônicos ligand-dependentes} (também chamados de ionotrópicos) são proteínas
multiméricas que formam um poro condutor de íons. A ligação do ligante --- um
neurotransmissor como acetilcolina (receptor nicotínico), GABA (receptor
GABA$_A$) ou glutamato (receptores NMDA e AMPA) --- causa uma mudança
conformacional que abre o canal em milissegundos, gerando fluxo de íons e
alterando o potencial de membrana. São os receptores de maior velocidade
no sistema nervoso.

Os \textbf{receptores nucleares}, por outro lado, são fatores de transcrição
ativados por ligantes hidrofóbicos (esteroides, hormônios tireoidianos,
retinoides) que atravessam a membrana plasmática por difusão. O ligante se une
ao receptor no citoplasma ou no núcleo, induzindo dimerização e recrutamento de
coativadores. O complexo migra para \emph{hormone response elements} (HREs) no
DNA, modulando a transcrição gênica --- uma resposta que leva horas, mas é
duradoura. Exemplos incluem o receptor de glicocorticoides (GR, alvo de
fármacos antiinflamatórios), o receptor de estrógeno (ER, alvo do tamoxifeno em
câncer de mama) e o PPAR$\gamma$ (regulação de metabolismo lipídico, alvo de
tiazolidinedionas usadas no diabetes tipo 2).

\section{Redes de Transdução de Sinal}
\label{sec:09-transducao}

Os sinais extracelulares raramente são transduzidos por um único caminho linear.
A célula opera com uma rede integrada de vias que se entrecruzam, amplificam e
modulam mutuamente --- o chamado \textbf{crosstalk}. As principais vias
canônicas --- MAPK/ERK, PI3K/Akt, JAK-STAT, Wnt, Notch, TGF-$\beta$/Smad, NF-$\kappa$B
e Hedgehog --- formam o esqueleto da resposta celular à maioria dos estímulos
extracelulares.

\subsection{Via MAPK/ERK: Proliferação e Diferenciação}
\label{subsec:09-mapk}

A cascata MAPK (proteína quinase ativada por mitógeno) é o protótipo de cascata
de fosforilação em biologia. A sequência canônica é:

$$\mathrm{RTK} \xrightarrow{\text{Grb2/SOS}} \mathrm{Ras\text{-}GTP}
\xrightarrow{} \mathrm{Raf\;(MAPKKK)} \xrightarrow{P}
\mathrm{MEK\;(MAPKK)} \xrightarrow{P\,P}
\mathrm{ERK\;(MAPK)} \xrightarrow{P}
\mathrm{fatores\;de\;transcrição}$$

Em detalhe: o receptor ativado recruta a adaptadora Grb2 e o fator de troca SOS;
SOS carrega a pequena GTPase Ras com GTP, ativando-a; Ras-GTP recruta Raf no
lado citoplasmático da membrana; Raf fosforila MEK em duas serinas; MEK
duplamente fosforilada fosforila ERK em um resíduo de treonina e outro de
tirosina (motivo TXY); ERK ativada dimeriza e transloca para o núcleo, onde
fosforila fatores de transcrição como Elk-1 e c-Fos, disparando um programa de
expressão gênica voltado para proliferação e diferenciação.

A via MAPK está entre as mais frequentemente desreguladas em câncer. Mutações
em RAS ocorrem em cerca de 30\% dos tumores humanos (KRAS em câncer de pâncreas;
NRAS em melanoma); mutações em RAF são encontradas em mais de 60\% dos melanomas;
mutações em MEK e ERK são raras mas existem. BRAF V600E é o exemplo
pedagógico: uma única transversão T $\to$ A no nucleotídeo 1799 produz uma
substituição V600E que mimetiza a fosforilação do segmento de ativação, mantendo
a quinase em estado constitutivamente ativo. O vemurafenib --- inibidor seletivo
de BRAF V600E --- exemplifica o paradigma da medicina de precisão: o fármaco é
eficaz apenas em tumores que carregam a mutação.

\subsection{Via PI3K/Akt: Sobrevivência e Metabolismo}
\label{subsec:09-pi3k}

A via PI3K/Akt controla simultaneamente sobrevivência celular, síntese
proteica, metabolismo de glicogênio e captação de glicose. A cascata é
iniciada por RTKs ou GPCRs que recrutam PI3K (fosfatidilinositol 3-quinase),
a qual fosforila o fosfolipídeo de membrana PIP$_2$ para gerar PIP$_3$. PIP$_3$
serve como plataforma citoplasmática que recruta PDK1 e Akt (também chamada
de PKB) via seus domínios PH. PDK1 fosforila Akt em T308; mTORC2 contribui
com fosforilação em S473 para ativação completa.

Entre os alvos de Akt estão mTOR (síntese proteica, via S6K e 4E-BP), GSK3
(glicogênio sintase, gliconeogênese), Bad (inibidor de apoptose), FoxO
(fator de transcrição pró-apoptótico) e MDM2 (regulador negativo de p53).
A desregulação é oncogênica: \emph{PTEN}, a fosfatase que reverte PIP$_3$ em
PIP$_2$, é um dos supressores tumorais mais frequentemente mutados em câncer.
A perda de PTEN mantém Akt constitutivamente ativa, conferindo às células
tumorais resistência à apoptose e vantagem proliferativa.

\subsection{Via JAK-STAT: Sinalização de Citocinas}
\label{subsec:09-jakstat}

A via JAK-STAT é a via de transdução mais curta e direta conhecida: receptor
$\to$ JAK $\to$ STAT $\to$ núcleo. Receptores de citocinas (interferons,
interleucinas, hormônios como eritropoietina e hormônio do crescimento) não
possuem atividade quinase intrínseca, mas estão associados constitutivamente a
quinases da família JAK (Janus quinases). A ligação da citocina dimeriza o
receptor, aproximando os JAKs associados e permitindo fosforilação cruzada.
Os JAKs então fosforilam resíduos de tirosina na cauda citoplasmática do
receptor, criando sítios de ancoragem para proteínas STAT (signal transducers
and activators of transcription) via seus domínios SH2. STATs recrutadas são
fosforiladas por JAK, dimerizam, e migram para o núcleo onde ativam genes
alvo com sítios GAS (interferon-\textit{gamma}-activated sequences).

\subsection{Vias Wnt, Notch e TGF-$\beta$/Smad}
\label{subsec:09-wtnnotch}

Três vias adicionais merecem destaque por sua importância em desenvolvimento
e em patologia:

A \textbf{via Wnt/$\beta$-catenina} controla a biogênese e estabilidade do
núcleo coativador $\beta$-catenina. Na ausência de Wnt, $\beta$-catenina é
seqüestrada por um complexo de destruição --- formado por APC, axina e GSK3---
que a fosforila, marca com ubiquitina e a direciona ao proteassomo. Com Wnt
ligado ao receptor Frizzled e correceptor LRP5/6, a proteína Dishevelled
inibe o complexo de destruição; $\beta$-catenina acumula no citoplasma,
transloca para o núcleo e, em parceria com fatores TCF/LEF, ativa genes
alvo. Mutações em APC estão presentes em mais de 80\% dos cânceres
colorretais.

A \textbf{via Notch} tem arquitetura única: o receptor é clivado ao receber o
ligante Delta (na membrana da célula vizinha). A porção intracelular (NICD)
se liberta, migra para o núcleo e associa-se ao fator CSL para ativar
transcrição. Não há cascata intermediária --- o sinal chega direto ao núcleo.
Notch media decisões de destino celular: célula-filha adota identidade A ou B
conforme receba mais ou menos sinalização Notch.

A \textbf{via TGF-$\beta$/Smad} usa a família de fatores Smad como transdutores.
TGF-$\beta$ liga-se a receptor serina/treonina quinase tipo II, que recruta e
fosforila receptor tipo I; este fosforila R-Smads (Smad2/3), que se associam à
Co-Smad (Smad4) e migram para o núcleo. I-Smads (Smad6/7) servem como
inibidores por feedback. TGF-$\beta$ tem papel dual em câncer: supressor
tumoral em estágios iniciais, mas promotor de metástase e EMT em tumores
avançados.

\begin{exemploenv}[Wnt em câncer colorretal]
Mais de 80\% dos cânceres colorretais abrigam mutações em APC que tornam o
complexo de destruição disfuncional. $\beta$-catenina acumula-se constitutivamente,
dirige a célula para um programa proliferativo sostenido e é suficiente para
iniciar a tumorigênese intestinal em modelos animais.
\end{exemploenv}


\section{Cascatas de Proteínas Quinases}
\label{sec:09-cascata-quinases}

\begin{definicaoenv}[Cascata de Quinases]
Sequência de reações bioquímicas na qual uma quinase ativa a seguinte por
fosforilação, resultando em amplificação progressiva do sinal. A arquitetura
em camadas é o que permite que milhares de eventos moleculares sejam
produzidos a partir de poucas ligações ao receptor.
\end{definicaoenv}

As cascatas de MAPKs são o exemplo canônico. Em cada nível, uma MAPKKK
ativa uma MAPKK, que por sua vez ativa uma MAPK. A arquitetura garante três
propriedades fundamentais: \textbf{amplificação} de sinal (uma molécula
ativada produz dezenas a centenas de cópias ativas no nível seguinte),
\textbf{ultrassensibilidade} (resposta sigmoidal próxima a um degrau), e
\textbf{integração} (a cascata pode ser modulada em cada nó). O componente
básico é uma reação de fosforilação/desfosforilação catalisada por uma
quinase e revertida por uma fosfatase.

\subsection{Ultrassensibilidade: Modelo de Goldbeter-Koshland}
\label{subsec:09-ultrassensibilidade}

Goldbeter e Koshland mostraram em 1981 que um simples ciclo de
fosforilação/desfosforilação pode exibir comportamento ultra-sensível ---
uma resposta sigmoidal praticamente um degrau --- quando as constantes de
Michaelis $K_1$ e $K_2$ são pequenas em relação à concentração total do
substrato. Para um substrato $X$ que é fosforilado e depois desfosforilado,
a fração ativa no equilíbrio é:

\begin{equation}
\frac{[X^*]}{[X]_{\mathrm{total}}} = \frac{2 v q}{v + q - \sqrt{(v - q)^2
                              + 4 v q \left(\frac{[X]_{\mathrm{tot}}}{K_1}\right)
                              \left(\frac{[X]_{\mathrm{tot}}}{K_2}\right)^{-1}}}
\end{equation}

onde $v$ e $q$ são as velocidades (opostas) de fosforilação e
desfosforilação. O resultado central é que o \emph{coeficiente de Hill efetivo},
medido experimentalmente como a inclinação da curva dose-resposta em escala
logarítmica, pode exceder em muito a estequiometria (1):

\begin{equation}
n_H \;=\; 1 \;+\; \frac{\ln 2}{\ln\!\bigl(1 + [X]_{\mathrm{tot}}/K_1\bigr)
                              \;+\; \ln\!\bigl(1 + [X]_{\mathrm{tot}}/K_2\bigr)}
\end{equation}

Quando $[X]_{\mathrm{tot}} \gg K_1, K_2$, a resposta torna-se
quase-digital: a transição de OFF para ON ocorre em uma estreita janela de
concentração do estímulo. Na cascata MAPK, esse efeito multiplica-se: três
níveis em série podem produzir $n_H$ efetivo de 5--10.

\subsection{Biestabilidade e Histerese}
\label{subfigure:09-biestabilidade}

Quando a ultra-sensibilidade é combinada a um \textbf{feedback positivo} --- em
que o produto da cascata também ativa seu próprio precursor --- a resposta pode
se tornar biestável: existem dois estados estáveis (OFF e ON) separados por um
estado instável intermediário. A célula então exibe \textbf{histerese}: o
limiar para ligar o sistema é maior do que o limiar para desligá-lo.

\begin{definicaoenv}[Biestabilidade]
Sistema com dois estados estáveis (ON e OFF) separados por um ponto de
instabilidade. Pequenas flutuações podem produzir comutação discreta entre os
dois estados; uma vez ligados, eles persistem mesmo quando o estímulo
decai abaixo do limiar original de ativação, conferindo memória celular.
\end{definicaoenv}

Os requisitos para biestabilidade são ultra-sensibilidade ($n_H > 1$),
feedback positivo e saturação de fosfatases. Consequência funcional:
decisões de longo prazo se tornam possíveis. A resposta imune adaptativa
envolve biestabilidades em populações de linfócitos T; o sistema de check-points
do ciclo celular é biestável; muitos processos de desenvolvimento
operam como comutadores biestáveis que precisam ser travados, não sensíveis a
estímulos transitórios.

\subsection{Fosforilação Distributiva vs. Processiva}
\label{sub:image9-distributiva}

Há duas maneiras arquiteturais pelas quais uma quinase pode fosforilar um
substrato que tem dois sítios de modificação. Na fosforilação \textbf{distributiva}, o substrato se dissocia da quinase após cada fosforilação; as duas fosforilações
são eventos independentes. Na fosforilação \textbf{processiva}, o substrato
permanece ligado à quinase até que ambos os sítios estejam fosforilados.

A arquitetura distributiva gera automaticamente ultra-sensibilidade: a
probabilidade de ambas as fosforilações ocorrerem é o quadrado da probabilidade
de cada uma, e a curva resultante tem a inclinação sigmoidal acentuada que
descrevemos. A cascata MAPK é distributiva --- MEK se dissocia de ERK após
cada fosforilação --- e este é o motivo pelo qual a MAPK exibe $n_H$ efetivo
tão elevado. Por contraste, fosforilação processiva gera resposta mais gradual
(mais próxima de Michaelis-Menten, $n_H \approx 1$).

\subsection{Simulação Computacional}
\label{sub:image9-simulacao}

O modelo simplificado abaixo de três níveis de cascata ilustra como reproduzir
as propriedades estudadas:

\begin{lstlisting}[language=Python, caption=Cascata MAPK de 3 niveis]
import numpy as np
from scipy.integrate import odeint
import matplotlib.pyplot as plt

def mapk_cascade(y, t, stimulus):
    K1, K1p, K2, K2p, K3, K3p = y
    # Ativacao (k1 depende do estimulo), inativacao linear
    v_in_1  = 0.1 * stimulus * K1
    v_out_1 = 0.1 * K1p
    v_in_2  = 0.5 * K1p * K2
    v_out_2 = 0.2 * K2p
    v_in_3  = 1.0 * K2p * K3
    v_out_3 = 0.3 * K3p
    return [-v_in_1 + v_out_1,
             v_in_1 - v_out_1 - v_in_2 + v_out_2,
            -v_in_2 + v_out_2,
             v_in_2 - v_out_2 - v_in_3 + v_out_3,
            -v_in_3 + v_out_3,
             v_in_3 - v_out_3]

y0 = [100, 0, 100, 0, 100, 0]
t = np.linspace(0, 50, 500)
stim = 1.0

sol = odeint(mapk_cascade, y0, t, args=(stim,))
plt.figure(figsize=(11, 5))

# Painel 1: dinamica temporal
plt.subplot(1, 2, 1)
plt.plot(t, sol[:, 1], 'b-',  lw=2, label='K1* (MAPKKK)')
plt.plot(t, sol[:, 3], 'g-',  lw=2, label='K2* (MAPKK)')
plt.plot(t, sol[:, 5], 'r-',  lw=2, label='K3* (MAPK)')
plt.xlabel('Tempo')
plt.ylabel('Concentracao (forma ativa)')
plt.title('Dinamica da cascata MAPK')
plt.legend(); plt.grid(alpha=0.3)

# Painel 2: curva dose-resposta
plt.subplot(1, 2, 2)
stimuli = np.logspace(-2, 1, 25)
resp = [odeint(mapk_cascade, y0, t, args=(s,))[-1, 5] for s in stimuli]
plt.semilogx(stimuli, resp, 'ro-', lw=2)
plt.xlabel('Estimulo (log)')
plt.ylabel('MAPK ativa no equilibrio')
plt.title('Ultrasensibilidade da cascata')
plt.grid(alpha=0.3)
plt.tight_layout()
plt.savefig('cascata_mapk.png', dpi=150)
plt.show()
\end{lstlisting}

O gráfico gerado mostra que a curva dose-resposta da cascata MAPK completa é
substancialmente mais sigmoidal do que a de um único estágio isolado: a
cascata amplifica o efeito não-linear de cada estágio.


\section{Segundos Mensageiros}
\label{sec:09-segundos-mensageiros}

\begin{definicaoenv}[Segundo Mensageiro]
Molécula intracelular pequena e difusível, produzida em resposta à ativação
de um receptor, que propaga e amplifica o sinal a efetores citoplasmáticos.
Permite que um sinal de membrana coordené respostas em todo o volume celular.
\end{definicaoenv}

Os principais segundos mensageiros são cAMP (AMP cíclico), cGMP (GMP cíclico),
Ca$^{2+}$, IP$_3$ (inositol 1,4,5-trifosfato), DAG (diacilglicerol) e NO (óxido
nítrico). Cada um possui um sistema próprio de síntese, degradação e alvos,
mas todos compartilham três propriedades: amplificação, difusão e integração.

\subsection{cAMP e a Via da Adenilil Ciclase}
\label{subfigure:09-camp}

A produção de cAMP é desencadeada por GPCRs acoplados a $G_s$: a subunidade
$G_s\alpha$-GTP ativa adenilil ciclase (AC), que catalisa a ciclização de ATP
a cAMP. cAMP liga-se às subunidades regulatórias da proteína quinase A (PKA),
liberando as subunidades catalíticas ativas. As subunidades catalíticas
fosforilam alvos como o fator de transcrição CREB (que ativa genes
responsivos a cAMP), enzimas metabólicas (incluindo a fosforilação que
inativa a glicogênio sintase) e canais iônicos. A terminação é mediada por
fosfodiesterases (PDEs) que degradam cAMP em AMP, encerrando o sinal.

A via do cAMP é central para muitas funções hormonais: a epinefrina ativa
$G_s$ nos hepatócitos, mobilizando glicose por glicogenólise; o glucagon usa o
mesmo mecanismo nos estados de jejum; o ACTH estimula a síntese de cortisol
nas adrenais via cAMP. A toxina da cólera é justamente uma toxina que
ADP-ribosila $G_s\alpha$, bloqueando sua GTPase --- o resultado é uma
produção descontrolada de cAMP que provoca a diarreia maciça característica
da doença.

\subsection{Sinalização por Ca$^{2+}$: Oscilações e Ondas}
\label{subfigure:09-calcio}

O íon Ca$^{2+}$ tem papel duplo como mensageiro e como cofator enzimático. No
estado de repouso, a concentração citosólica de Ca$^{2+}$ é mantida em torno de
100 nM, três a quatro ordens de grandeza abaixo do cálcio extracelular (1--2 mM)
e do estoque no retículo endoplasmático (0.5 mM). Essa diferença de
concentração gera o potencial eletroquímico que dirige a abertura de canais
de Ca$^{2+}$, gerando sinais transitórios ou oscilatórios.

A produção de IP$_3$ por fosfolipase C (PLC) é o principal gatilho. IP$_3$
difunde do interior da membrana plasmática para o retículo, liga-se a
receptores IP$_3$R (canais de Ca$^{2+}$ regulados por ligante) e libera Ca$^{2+}$
diretamente no citoplasma. O Ca$^{2+}$ liberado pode então ativar a PLC
local, produzindo mais IP$_3$ --- esse \textbf{mecanismo de liberação
autoinduzida de Ca$^{2+}$} (CICR, \emph{calcium-induced calcium release}) é
responsável pelas oscilações observadas em muitos tipos celulares.

As oscilações de Ca$^{2+}$ funcionam como código de frequência: baixa frequência
codifica estímulo fraco, alta frequência codifica estímulo forte. Dolmetsch
et al.\ (1998) mostraram que células de HEK respondem à histamina com trens de
\emph{spikes} de Ca$^{2+}$ cuja frequência (e não amplitude) determina a
especificidade transcricional: frequência baixa ativa NF-$\kappa$B,
frequência alta ativa NFAT. A amplitude do sinal permanece saturada; é o
ritmo que contém a informação.

Um modelo simplificado de oscilações inclui três termos: influxo basal do
meio extracelular, efluxo basal, liberação mediada por IP$_3$ e CICR
(proporcional a $[\mathrm{Ca}^{2+}]^2$), e recaptação pelo retículo (proporcional
linearmente a $[\mathrm{Ca}^{2+}]$). As equações:
\begin{align}
\frac{d[\mathrm{Ca}^{2+}]}{dt} &= v_{\mathrm{in}} - v_{\mathrm{out}}
+ v_{\mathrm{release}} - v_{\mathrm{uptake}} \\
v_{\mathrm{release}} &= k_{\mathrm{rel}} \cdot [\mathrm{IP}_3] \cdot
\frac{[\mathrm{Ca}^{2+}]^2}{K_1^2 + [\mathrm{Ca}^{2+}]^2} \\
v_{\mathrm{uptake}} &= k_{\mathrm{uptake}} \cdot [\mathrm{Ca}^{2+}]
\end{align}

geram oscilações sustentadas para uma faixa de valores de $[\mathrm{IP}_3]$.
A simulação computacional abaixo ilustra como a frequência aumenta com
$[\mathrm{IP}_3]$:

\begin{lstlisting}[language=Python, caption=Oscilacoes de Ca2+ via CICR]
import numpy as np
from scipy.integrate import odeint
import matplotlib.pyplot as plt

def ca_oscillations(y, t, IP3):
    Ca = y[0]
    k_rel, k_uptake, K1, v_in, v_out = 2.0, 0.5, 0.3, 0.1, 0.05
    v_rel    = k_rel * IP3 * Ca**2 / (K1**2 + Ca**2)
    v_uptake = k_uptake * Ca
    return [v_in - v_out + v_rel - v_uptake]

t = np.linspace(0, 200, 2000)
fig, axes = plt.subplots(1, 3, figsize=(15, 4))
for i, IP3 in enumerate([0.5, 1.0, 1.5]):
    sol = odeint(ca_oscillations, [0.1], t, args=(IP3,))
    axes[i].plot(t, sol[:, 0], 'b-', lw=1.5)
    axes[i].set_xlabel('Tempo')
    axes[i].set_ylabel('[Ca2+]')
    axes[i].set_title(f'IP3 = {IP3}')
    axes[i].grid(alpha=0.3)
plt.tight_layout()
plt.savefig('oscilacoes_ca.png', dpi=150)
plt.show()
\end{lstlisting}

\subsection{Via IP$_3$/DAG e Ativação de PKC}
\label{subfigure:09-ip3-dag}

A hidrólise de PIP$_2$ por PLC gera duas moléculas com papéis distintos. O
IP$_3$ é hidrofílico e solúvel: difunde no citoplasma até o retículo
endoplasmático e libera Ca$^{2+}$ via receptores IP$_3$R, conforme descrito
acima. O DAG permanece na membrana: serve como âncora para a proteína
quinase C (PKC), recrutando-a do citoplasma para a membrana em sinapse
com o Ca$^{2+}$ recém-liberado. A combinação Ca$^{2+}$ + DAG ativa a PKC,
que fosforila substratos como MARCKS, proteínas do citoesqueleto e
fatores de transcrição. É a convergência dos dois ramos do PIP$_2$ que
garante especificidade: nem DAG sozinho, nem Ca$^{2+}$ sozinho, alcançam
ativação robusta de PKC.

\subsection{cGMP e Sinalização por Óxido Nítrico}
\label{subfigure:09-cgmp}

A via do cGMP difere da via do cAMP por seu segundo mensageiro ser um gás.
O óxido nítrico (NO) é produzido por NO sintases (NOS neuronal, endotelial
e induzível) a partir do aminoácido L-arginina. Por ser uma molécula
pequena, apolar e com meia-vida de poucos segundos, NO difunde rapidamente
através de membranas e atua em células vizinhas antes de ser oxidado a
nitrato. NO ativa a guanilil ciclase solúvel (sGC, contendo o grupo heme
que liga NO), que converte GTP em cGMP. cGMP ativa PKG (proteína quinase G),
que fosforila alvos como fosfolambdam (músculo cardíaco), canais iônicos
e proteínas do citoesqueleto.

A relevância clínica é enorme. NO produzido pelo endotélio vascular é o
principal mediador de vasodilatação: o nitroprussiato de sódio e a nitroglicerina
são doadores de NO; os inibidores de PDE5 (sildenafil/Viagra, tadalafil/Cialis)
bloqueiam a degradação de cGMP, prolongando o relaxamento vascular e tratando
disfunção erétil e hipertensão pulmonar. NO tem também papel na memória
(sinapses hipocampais), na resposta imune (macrófagos produzem NO para matar
bactérias) e na neurotransmissão não-adrenérgica não-colinérgica (NANC).


\section{Dinâmica e Propriedades Emergentes}
\label{sec:09-dinamica}

As propriedades mais notáveis das redes de sinalização --- oscilações,
adaptação, biestabilidade, robustez a ruído, memória --- não são inscritas em
nenhum componente individual. Emergem da interação entre componentes e da
arquitetura da rede. Esta seção descreve os principais padrões dinâmicos e as
condições para sua ocorrência.

\subsection{Loops de Feedback Positivo e Negativo}
\label{sub:image9-feedback}

O destino dinâmico de uma via de sinalização é em larga medida ditado pela
composição de seus feedbacks. \textbf{Feedbacks negativos} estabilizam o
sistema: reduzem a variabilidade em torno do estado estacionário, promovem
homeostase e adaptação. Exemplos incluem a inibição por produto (o produto
final reduz sua própria síntese), a desensibilização de receptores (GPCRs
fosforilados por GRKs perdem afinidade por proteínas G) e os I-Smads na
via TGF-$\beta$. São predominantes em sistemas homeostáticos: detecção de
glicose, regulação de volume celular, termorregulação.

\textbf{Feedbacks positivos} amplificam o sinal inicial e desestabilizam o
estado de equilíbrio. CICR no Ca$^{2+}$, caspases ativadas por caspases na
apoptose e os checkpoints do ciclo celular (que precisam comutar ON/OFF com
decisão persistente) dependem de feedback positivo. A combinação de ambos os
tipos de feedback em uma mesma rede é o que sustenta oscilações: o feedback
negativo com atraso temporal gera ciclo limite, enquanto o positivo garante
que a amplitude seja grande o suficiente para produzir um sinal detectável.

\subsection{Osciladores Biológicos}
\label{subfigure:09-oscilacoes}

Vários sistemas de sinalização exibem comportamento oscilatório espontâneo
mesmo na presença de estímulo sustentado. Os osciladores mais estudados são:

\begin{itemize}
\item \textbf{p53--Mdm2:} dano ao DNA gera uma série de pulsos de p53, com período
de algumas horas; p53 transativa \emph{MDM2}, que ubiquitiniza p53 e a marca
para degradação; com a queda de p53, MDM2 é desativada; p53 acumula e produz
outro pulso. Pesquisas de Lahav e colaboradores (2004) mostraram que o número
de pulsos e não a amplitude é o que codifica a severidade do dano.
\item \textbf{NF-$\kappa$B:} TNF-$\alpha$ gera oscilações de localização nuclear
de NF-$\kappa$B com período de 100 min; a oscilação depende de múltiplos
loops de feedback negativo via I$\kappa$B$\alpha$ e A20.
\item \textbf{Hes1:} o oscilador de segmentação embrionária usa repressão
transcricional com atraso --- Hes1 reprime sua própria expressão, de forma
que a proteína Hes1 oscila em torno de 2 horas.
\item \textbf{ERK:} respostas a EGF tendem a ser transitórias; respostas a
NGF ou PDGF podem ser sustentadas; a duração da ativação determina se a
célula prolifera ou diferencia.
\end{itemize}

\subsection{Adaptação e Desensibilização}
\label{sub:image9-adaptacao}

A adaptação é a capacidade que o sistema tem de retornar ao estado basal
mesmo na presença de estímulo sustentado. Em sensoriamento, adaptação
permite perceber \emph{mudanças} em vez de intensidades absolutas --- é o
princípio dos nossos sentidos de tato (terminações nervosas detectam
mudanças de temperatura, não a temperatura em si) e de visão (fotorreceptores
adaptam-se à luz ambiente para detectar contrastes).

Os mecanismos moleculares são variados: fosforilação do receptor por GRKs
(receptores fosforilados recrutam $\beta$-arrestinas, que bloqueiam o
acoplamento a proteínas G); internalização de receptores por endocitose;
indução de fosfatases e inibidores por feedback transcricional; depleção de
estoques intracelulares de precursores (por exemplo, esgotamento de PIP$_2$
local ao GPCR). A escala temporal vai de segundos (fosforilação por GRKs) a
horas (indução transcricional).

\subsection{Morfógenos e Gradientes Espaciais}
\label{sub:image9-morfogenos}

\begin{definicaoenv}[Morfógeno]
Morfógeno é uma molécula sinalizadora que forma um gradiente de concentração
no tecido e especifica destinos celulares distintos em função da concentração
local. O mesmo sinal é lido por diferentes genes-alvo em diferentes limiares,
gerando padrões espaciais discretos a partir de distribuições contínuas.
\end{definicaoenv}

A lógica é simples: a célula detecta a concentração absoluta do morfógeno,
compara com limiares internos, e adota um destino transcricional compatível
com a faixa detectada. Bicoid em \emph{Drosophila} gera o eixo ântero-posterior
por um gradiente de concentração de proteína morfogenética, com tradução
localizada a partir de mRNA materno depositado no polo anterior. Sonic
Hedgehog (Shh) organiza o tubo neural: concentrações altas geram neurônios
motores ventrais; concentrações intermediárias, interneurônios; concentrações
baixas, neurônios dorsais. BMP (bone morphogenetic protein) gera o eixo
dorso-ventral em vários sistemas.

O modelo matemático típico é a equação de difusão-destruição em estado
estacionário:
\begin{equation}
D \nabla^2 M - k [M] = 0
\end{equation}

cuja solução unidimensional é uma exponencial decrescente
$M(x) = M_0 e^{-x/\sqrt{D/k}}$. O comprimento característico
$\lambda = \sqrt{D/k}$ define a distância em que o morfógeno cai a $1/e$
de seu valor na fonte --- é o parâmetro que limita o tamanho do tecido que
pode ser padronizado por esse mecanismo.

\subsection{Robustez e Sensibilidade}
\label{sub:image9-robustez}

Um sistema biológico precisa simultaneamente ser sensível a sinais
específicos e robusto a perturbações aleatórias. A engenharia dos sistemas
de sinalização satisfaz essa dupla exigência por meio de quatro mecanismos:

\begin{enumerate}
\item \textbf{Feedback negativo}: mantém concentrações próximas do
valor desejado, amortecendo flutuações.
\item \textbf{Redundância}: vias paralelas (PI3K/Akt vs MAPK para EGFR; NF-$\kappa$B
e JNK para TNF-$\alpha$) garantem que a falha de uma via não abole a resposta.
\item \textbf{Compartimentalização}: localização subcelular (microdomínios de
membrana, endossomos, sinapses) separa espacialmente cascatas que poderiam
interferir entre si.
\item \textbf{Buffering}: chaperonas e proteínas de suporte (Hsp90, 14-3-3) ligam
quinases em estados inativos, estabilizando a rede contra ativação espúria.
\end{enumerate}

A relação de compromisso entre robustez e sensibilidade tem paralelo com a
teoria de controle: sistemas robustos a ruído tendem a ser lentos; sistemas
rápidos tendem a ser instáveis. As redes de sinalização biológicas são
engenharia fina em torno desse compromisso, usando ultra-sensibilidade,
bistabilidade e oscilações para obter \emph{robustez dinâmica} --- respostas
rápidas a sinais específicos, mas com rejeição firme a flutuações espúrias.
Bhalla e Iyengar (1999) mostraram em artigo seminal como combinar módulos de
sinalização simples (cada um com sua resposta característica) gera
comportamentos globais complexos --- biestabilidade, oscilações, adaptação ---
que nenhum módulo isolado conseguiria.


\section{Exercícios}
\label{sec:09-exercicios}

\begin{exercicio}
\textbf{Ligação receptor-ligante.} Um receptor apresenta $K_d = 5$ nM e
$B_{\max} = 1000$ receptores por célula. (a) Calcule a fração de ocupação
quando $[L] = 5$ nM. (b) Determine a concentração de ligante necessária para
atingir 90\% de ocupação. (c) Se $k_{\mathrm{on}} = 10^6\;\mathrm{M}^{-1}\mathrm{s}^{-1}$,
determine $k_{\mathrm{off}}$ e comente seu significado funcional.
\end{exercicio}

\begin{exercicio}
\textbf{Amplificação em cascata.} Considere uma cascata MAPK com três níveis
(MAPKKK, MAPKK, MAPK) na qual cada quinase ativada fosforila em média 10
moléculas do nível seguinte. Se 5 moléculas de MAPKKK forem ativadas, estime
o número de moléculas de MAPKK e MAPK ativas, e o fator de amplificação total.
Discuta como cascatas com mais níveis ou ganhos maiores modificariam a resposta.
\end{exercicio}

\begin{exercicio}
\textbf{Ultrassensibilidade de Goldbeter-Koshland.} Considere um sistema
simplificado de fosforilação/desfosforilação com $[X]_{\mathrm{tot}} = 100$ nM,
$K_1 = K_2 = 10$ nM. (a) Plote a fração ativa em função da razão $v/q$
para $[S]$ variável. (b) Meça o coeficiente de Hill efetivo em escala
logarítmica. (c) Compare o resultado com o caso $[X]_{\mathrm{tot}} = 1$ nM
(mesmas $K_1, K_2$) e interprete as diferenças.
\end{exercicio}

\begin{exercicio}
\textbf{cAMP em estado estacionário.} Em cultura de hepatócitos, cAMP é
produzido por adenilil ciclase a $v_{\mathrm{synth}} = 0{,}5\;\mu\mathrm{M/s}$
e degradado por fosfodiesterase seguindo cinética de primeira ordem com
$k_{\mathrm{deg}} = 0{,}1\;\mathrm{s}^{-1}$. (a) Escreva a equação
diferencial para $[cAMP]$. (b) Calcule a concentração de estado
estacionário. (c) Determine a constante de tempo para atingir 90\% do
equilíbrio após estimulação em degrau. (d) Como a constante de tempo se
compara com tempos típicos de resposta hepática (segundos a minutos)?
\end{exercicio}

\begin{exercicio}
\textbf{Simulação computacional da cascata MAPK.} Implemente o sistema de EDOs
apresentado no código deste capítulo. (a) Reproduza a curva dose-resposta
para 25 valores de estímulo distribuídos em escala logarítmica de
$10^{-2}$ a $10$. (b) Meça o coeficiente de Hill efetivo da curva por ajuste
logístico. (c) Investigue como $n_H$ varia quando o número de níveis da
cascata é alterado entre 1, 2 e 4. (d) Adicione feedback negativo
(estimulado pelo produto final) e discuta como a amplitude e a duração da
resposta se modificam.
\end{exercicio}

\begin{exercicio}
\textbf{Oscilações de Ca$^{2+}$.} Modifique o modelo de oscilações de
Ca$^{2+}$ apresentado no capítulo. (a) Varie sistematicamente $[\mathrm{IP}_3]$
em uma faixa logarítmica e meça a frequência de oscilação. (b) Investigue o
efeito de $k_{\mathrm{uptake}}$ sobre o regime oscilatório. (c) Adicione
um termo de difusão (modelo 1D com dez compartimentos) e mostre como se
formam ondas espaciais de Ca$^{2+}$. (d) Compare seu modelo com os dados
experimentais de Dolmetsch et al.\ (1998) sobre codificação por frequência.
\end{exercicio}

\begin{exercicio}
\textbf{Modelo simplificado de via Wnt/$\beta$-catenina.} Formule um sistema
de EDOs com três variáveis: Wnt extracelular ($W$), $\beta$-catenina
citosólica ($B_c$) e $\beta$-catenina nuclear ($B_n$). Sob Wnt, a taxa de
degradação de $B_c$ é reduzida (efeito sobre o complexo de destruição); sob
ausência de Wnt, $B_c$ é fosforilado e ubiquitinado. (a) Plote a resposta
em degrau (Wnt ON/OFF). (b) Investigue como variações nos parâmetros
imitam mutações em APC (degradação constitutivamente reduzida).
\end{exercicio}

\begin{exercicio}
\textbf{Projeto integrador: modelagem de uma via completa.} Escolha uma das
vias discutidas no capítulo (MAPK, PI3K/Akt, JAK-STAT, Wnt) e construa um
modelo matemático com 6--12 componentes, com base em revisão da literatura.
(a) Implemente o modelo em Python usando \texttt{scipy.integrate.odeint}.
(b) Simule perturbações genéticas e farmacológicas: nocaute de uma
componente constitutivamente inativado, inibição parcial de quinase,
indução de superexpressão. (c) Analise a sensibilidade dos parâmetros
usando análise de variância global (PRCC) ou amostragem de Latin Hypercube.
(d) Gere diagnósticos quantitativos comparando previsões a dados
experimentais de bases públicas (GEO, PathwayCommons).
\end{exercicio}

\section{Notas Bibliográficas}
\label{sec:09-bibliografia}

Os fundamentos da sinalização celular são apresentados com profundidade nas
obras clássicas de Alberts et al.\ (\textit{Molecular Biology of the Cell},
7ª edição) e Lodish et al.\ (\textit{Molecular Cell Biology}, 9ª edição),
que permanecem referências obrigatórias para a visão integrada. A leitura
de \textit{An Introduction to Systems Biology} de Uri Alon (2ª edição) é
particularmente recomendada: oferece uma apresentação unificada dos princípios
de design de circuitos biológicos, com exemplos concretos em cascatas de
fosforilação, osciladores e biestabilidade. O livro \textit{Systems Biology: A
Textbook} de Klipp, Herrmann e colaboradores (2ª edição) é referência
indispensável para os aspectos quantitativos e computacionais.

Vários artigos seminais contribuíram para a compreensão moderna das cascatas
de sinalização como sistemas dinâmicos. O de Kholodenko (2006) em
\textit{Nature Reviews Molecular Cell Biology} é uma exposição influente das
propriedades espaço-temporais das redes de sinalização. A revisão clássica de
Ferrell (1996) sobre cascatas de MAPK como conversores de sinais graduais em
respostas em chave permanece leitura obrigatória. O artigo de Goldbeter e
Koshland (1981) em \textit{PNAS}, sobre ultra-sensibilidade em sistemas
bioquímicos, é o ponto de partida para qualquer análise matemática de
cascatas de quinases. O estudo computacional de Bhalla e Iyengar (1999) em
\textit{Science} demonstrou de forma pioneira como combinar módulos simples
de sinalização produz comportamentos complexos emergentes --- biestabilidade,
oscilações, adaptação --- que nenhum dos módulos isolados seria capaz de exibir.

Para a sinalização por Ca$^{2+}$, o trabalho de Dolmetsch et al.\ (1998) em
\textit{Nature} estabeleceu que a frequência das oscilações --- e não a amplitude ---
codifica informação transcricional. Os osciladores circadianos e de ciclo
celular são discutidos em profundidade por Tyson e Novak em \textit{The Cell
Cycle: Principles of Control} (Oxford). Para um tratamento computacional mais
aprofundado, \textit{Computational Cell Biology} de Fall, Marland, Wagner e
Tyson (Springer) é referência consolidada.

Recursos online valiosos incluem KEGG Pathway e Reactome, duas bases de
dados curadas manualmente que descrevem vias de sinalização humanas em
formato acessível por máquina; SignaLink, base dedicada à sinalização
multicamada em eucariotos; e PathwayCommons, agregador de vias de múltiplas
fontes. A integração dessas bases com modelos computacionais permite validar
predições teóricas em conjuntos de dados públicos.

\begin{thebibliography}{99}
\bibitem{alberts} Alberts B., Johnson A., Lewis J., Raff M., Roberts K.,
Walter P. (2022). \textit{Molecular Biology of the Cell}, 7ª ed.
Garland Science, Nova York.

\bibitem{lodish} Lodish H., Berk A., Kaiser C.A., Krieger M., Bretscher A.,
Ploegh H., Amon A., Martin K.C. (2021). \textit{Molecular Cell Biology},
9ª ed. W.H. Freeman, Nova York.

\bibitem{alon} Alon U. (2019). \textit{An Introduction to Systems Biology:
Design Principles of Biological Circuits}, 2ª ed. CRC Press, Boca Raton.

\bibitem{klipp} Klipp E., Herrmann H., Liebermeister W., Rother K., Gilles
E.D., Wierling C., Kowald A., Lehrach H. (2016). \textit{Systems Biology: A
Textbook}, 2ª ed. Wiley-VCH, Weinheim.

\bibitem{kholodenko} Kholodenko B.N. (2006). Cell-signalling dynamics in
time and space. \textit{Nature Reviews Molecular Cell Biology},
7:165--176.

\bibitem{ferrell} Ferrell J.E.\ Jr (1996). Tripping the switch fantastic:
how a protein kinase cascade can convert graded inputs into switch-like
outputs. \textit{Trends in Biochemical Sciences}, 21:460--466.

\bibitem{goldbeter} Goldbeter A., Koshland D.E.\ Jr (1981). Ultrasensitivity
in biochemical systems. \textit{Proceedings of the National Academy of
Sciences USA}, 78:6840--6844.

\bibitem{bhalla} Bhalla U.S., Iyengar R. (1999). Emergent properties of
networks of biological signaling pathways. \textit{Science}, 283:381--387.

\bibitem{dolmetsch} Dolmetsch R.E., Xu K., Lewis R.S. (1998). Calcium
oscillations increase the efficiency and specificity of gene expression.
\textit{Nature}, 392:933--936.

\bibitem{tyson} Tyson J.J., Novak B. (2020). \textit{The Cell Cycle:
Principles of Control}. Oxford University Press, Oxford.

\bibitem{fall} Fall C.P., Marland E.S., Wagner J.M., Tyson J.J. (2002).
\textit{Computational Cell Biology}. Springer, Nova York.

\bibitem{lahav} Lahav G., Rosenfeld N., Sigal A., Geva-Zatorsky N., Levine
A.J., Elowitz M.B., Alon U. (2004). Dynamics of the p53--Mdm2 feedback
loop in individual cells. \textit{Nature Genetics}, 36:147--150.

\bibitem{pulvermacher} KEGG Pathway Database. \textit{Kyoto Encyclopedia of
Genes and Genomes}. Disponível em: \url{https://www.genome.jp/kegg/pathway.html}.
Acessado em 2026.

\bibitem{reactome} Reactome Knowledgebase. Disponível em:
\url{https://reactome.org}. Acessado em 2026.

\bibitem{salink} SignaLink 2.0: a service to monitor signaling pathways.
Disponível em: \url{http://signalink.org}. Acessado em 2026.

\bibitem{pathcomm} Pathway Commons. Disponível em:
\url{https://www.pathwaycommons.org}. Acessado em 2026.
\end{thebibliography}
